% Bagian: propositional-logic
% Bab: propositional-logic
% Subbagian: preliminaries

\documentclass[../../../include/open-logic-section]{subfiles}

\begin{document}

\olfileid[id]{pl}{syn}{pre}

\olsection{Persiapan}

\begin{thm}[\emph{Prinsip induksi pada !!{formula}s}]
  \ollabel{thm:induction}  
  Jika suatu sifat~$P$ berlaku bagi semua !!{formula}s atomik dan memenuhi
  ketentuan berikut:
  \begin{enumerate}
    \tagitem{prvNot}{sifat itu berlaku bagi $\lnot !A$ setiap kali berlaku
      bagi~$!A$;}{}
    \tagitem{prvAnd}{sifat itu berlaku bagi $(!A \land !B)$ setiap kali
      berlaku bagi $!A$ dan~$!B$;}{}
    \tagitem{prvOr}{sifat itu berlaku bagi $(!A \lor !B)$ setiap kali
      berlaku bagi $!A$ dan~$!B$;}{}
    \tagitem{prvIf}{sifat itu berlaku bagi $(!A \lif !B)$ setiap kali
      berlaku bagi $!A$ dan~$!B$;}{}
    \tagitem{prvIff}{sifat itu berlaku bagi $(!A \liff !B)$ setiap kali
      berlaku bagi $!A$ dan~$!B$;}{}
  \end{enumerate}
  maka $P$ berlaku bagi semua !!{formula}s.
\end{thm}

\begin{proof}
  Misalkan $S$ ialah koleksi semua !!{formula}s yang memiliki sifat~$P$.
  Jelas bahwa $S \subseteq \Frm[L_0]$. $S$~memenuhi semua syarat dalam
  \olref[fml]{defn:formulas}: himpunan itu memuat semua !!{formula}s atomik dan
  tertutup terhadap !!{operator}s. $\Frm[L_0]$ merupakan kelas terkecil yang
  demikian, sehingga $\Frm[L_0] \subseteq S$. Jadi $\Frm[L_0] = S$, dan
  setiap formula memiliki sifat~$P$.
\end{proof}

\begin{prop}\ollabel{prop:balanced}
   Setiap !!{formula} dalam~$\Frm[L_0]$ bersifat \emph{seimbang}, dalam arti
   memiliki tanda kurung kiri sebanyak tanda kurung kanan.
\end{prop}

\begin{prob}
Buktikan \olref[pl][syn][pre]{prop:balanced}
\end{prob}

\begin{prop} \ollabel{prop:noinit}
Tidak ada segmen awal sejati dari !!a{formula} yang merupakan !!a{formula}.
\end{prop}

\begin{prob}
  Buktikan \olref[pl][syn][pre]{prop:noinit}
\end{prob}

\begin{prop}[Keterbacaan Unik]
Setiap !!{formula}~$!A$ dalam $ \Frm[L_0]$ memiliki tepat satu penguraian
sebagai salah satu bentuk berikut:
\begin{enumerate}
\tagitem{prvFalse}{$\lfalse$.}{}

\tagitem{prvTrue}{$\ltrue$.}{}

\item $\Obj p_n$ untuk suatu $\Obj p_n \in  \PVar$.
  
\tagitem{prvNot}{$\lnot !B$ untuk suatu !!{formula}~$!B$.}{}

\tagitem{prvAnd}{$(!B \land !C)$ untuk !!{formula}s $!B$ dan~$!C$ tertentu.}{}

\tagitem{prvOr}{$(!B \lor !C)$ untuk !!{formula}s $!B$ dan~$!C$ tertentu.}{}

\tagitem{prvIf}{$(!B \lif !C)$ untuk !!{formula}s $!B$ dan~$!C$ tertentu.}{}

\tagitem{prvIff}{$(!B \liff !C)$ untuk !!{formula}s $!B$ dan~$!C$ tertentu.}{}
\end{enumerate}
Selain itu, penguraian ini \emph{unik}.
\end{prop}

\begin{proof}
Melalui induksi pada $!A$. Sebagai contoh, andaikan $!A$ memiliki dua
penguraian berbeda sebagai $(!B \lif !C)$ dan $(!B' \lif !C')$. Maka $!B$
dan $!B'$ harus sama (jika tidak, salah satunya merupakan segmen awal sejati
dari yang lain); jadi, jika kedua penguraian $!A$ itu berbeda, penyebabnya
pastilah bahwa $!C$ dan $!C'$ merupakan penguraian berbeda dari barisan
simbol yang sama, dan hal ini mustahil menurut hipotesis induksi.
\end{proof}


\begin{defn}[Substitusi Seragam]
Jika $!A$ dan $!B$ merupakan !!{formula}s, dan $\Obj p_i$ merupakan suatu
!!{propositional
variable}, maka $\Subst{!A}{!B}{\Obj p_i}$ menyatakan hasil
penggantian setiap kemunculan $\Obj p_i$ dengan satu kemunculan $!B$ di
dalam $!A$; serupa dengan itu, substitusi serentak $\Obj p_1$, \dots,~$\Obj p_n$
dengan !!{formula}s $!B_1$, \dots,~$!B_n$ dinyatakan dengan
$\SSubst{!A}{\subst{!B_1}{\Obj p_1},\dots,\subst{!B_n}{\Obj p_n}}$.
\end{defn}

\begin{prob} Untuk masing-masing dari lima !!{formula}s di bawah ini,
 tentukan apakah !!{formula} tersebut dapat dinyatakan sebagai substitusi
 \( \Subst{!A}{!B}{\Obj p_i} \), dengan \( !A \) berupa (i)~\( \Obj p_0 \);
 (ii)~\( ( \lnot \Obj p_0 \land \Obj p_1) \); dan (iii)~\( ( ( \lnot \Obj
 p_0 \lif \Obj p_1 ) \land \Obj p_2 ) \). Dalam setiap kasus, tentukan
 substitusi yang bersangkutan.
  \begin{enumerate}
    \item \( \Obj p_1 \) 
    \item \( ( \lnot \Obj p_0 \land \Obj p_0 ) \)
    \item \( ( ( \Obj p_0 \lor \Obj p_1 ) \land \Obj p_2 )  \)
    \item \( \lnot ( ( \Obj p_0 \lif \Obj p_1 ) \land \Obj p_2 ) \)
    \item \( (( \lnot ( \Obj p_0 \lif \Obj p_1 ) \lif ( \Obj p_0 \lor \Obj p_1 )) \land \lnot ( \Obj p_0 \land \Obj p_1 )) \)
  \end{enumerate}
\end{prob}

\begin{prob}
Berikan definisi yang ketat secara matematis bagi $\Subst{!A}{!B}{p}$ melalui
induksi.
\end{prob}

\end{document}
